\begin{frame}{회귀 vs 분류: 같은 모델, 다른 출력}
  지도학습은 다시 출력이 \textbf{연속값이냐 범주냐}로 나뉜다:
  \begin{itemize}
    \item \kb{회귀(regression)}: \textbf{숫자}를 예측 (집값, 기온).
    \item \kb{분류(classification)}: \textbf{범주}를 예측 (스팸/정상,
      개/고양이/새).
  \end{itemize}
  \vskip0.4em
  같은 선형 모델 $w^\top x + b$ 라도, \textbf{출력을 그대로 쓰면 회귀},
  \textbf{시그모이드를 씌워 0$\sim$1 확률로 해석하면 분류} (둘 다
  Week 2). 앞으로 몇 장 동안 반복되는 패턴:
  \begin{block}{모델 구조는 비슷한데, 출력을 해석하는 방식과
  손실함수(loss function)만 문제에 맞게 바뀐다.}
  \end{block}
  \vskip0.4em
  \begin{itemize}
    \item 예: 학습된 $w = [0.5, -2.0]$, \; $b = 1.0$, \; 새 데이터
      $x = [3, 0.5]$.
      \[w^\top x + b = 0.5 \times 3 + (-2.0) \times 0.5 + 1.0
        = 1.5 - 1.0 + 1.0 = 1.5\]
  \end{itemize}
\end{frame}
